\section{Four Levels of ``The Output Is the Same''}
\label{sec:equivalence-levels}

The exact Quirk export analysis makes a crucial methodological point:
matching amplitude-derived bitstring distributions do not establish every
notion of equivalence. The advanced circuit satisfies the weakest level below
while failing the next one. We propose four levels.

\begin{table}[ht]
\centering
\caption{Four levels of output equivalence. Each level is stronger than the
one above it.}
\label{tab:equivalence-levels}
\begin{tabularx}{\textwidth}{@{}>{\bfseries}p{0.12\textwidth}p{0.25\textwidth}X@{}}
\toprule
Level & What is equal? & What can still differ?\\
\midrule
E0: Histogram & Selected measured bitstring probabilities & Relative phases,
coherence, hidden correlations, behavior in other bases\\
E1: Reduced state & Full density matrix of the legitimate output register &
Attacker correlations compatible with mixed/classical outputs; hidden registers\\
E2: Process & The legitimate channel on a specified input family, ideally with
entangled references & Behavior outside the tested family; implementation-level
side channels\\
E3: Execution & Global authorized semantics across gate, mapping, pulse, and
hardware layers & Only physical noise and explicitly accepted leakage\\
\bottomrule
\end{tabularx}
\end{table}

\subsection{E0 is especially dangerous}

For the attacked $\ket{+}$ example, the baseline and attacked states both yield
$0$ and $1$ with probability $1/2$ in the computational basis. Yet the baseline
has $\langle X\rangle=1$, while the attacked reduced state has
$\langle X\rangle=0$. The same histogram hides complete loss of phase
coherence.

\subsection{E1 does not always prove secrecy}

If the legitimate output is a mixed state encoding a classical distribution,
an attacker may share correlations while leaving the victim marginal unchanged.
No-broadcasting limits this for noncommuting state families
\cite{barnum1996broadcasting}, but classical components remain readable without
disturbance \cite{koashi2002operations}.

\subsection{E2 should include reference systems}

Testing only isolated inputs can miss failures on entangled inputs. The
entangled-reference proof in \cref{cor:reference-preservation} is therefore not
an ornamental detail; it is the appropriate standard for claiming that a
subroutine implements a quantum channel correctly.

\subsection{E3 is a cross-layer security goal}

Even a verified gate-level circuit can be implemented by altered routing or
pulse schedules. Recent work studies quantum Trojans, malicious compilation,
and mismatches between gate-level declarations and pulse-level execution
\cite{ashsaki2023primer,john2025trojan,xu2024pulseattacks}. A user who verifies
only source code or a final histogram has not verified the full execution.
